\section{Certificates and the certification gap}\label{sec:certificates}

Fragility has a robustness companion. The same operator-norm bracket that upper-bounds the network
Lipschitz constant gives, at every input, a radius inside which no perturbation can change the
output; this is a genuine certificate, valid for all inputs with no exceptional set. Comparing it with
the fragility radius of \Cref{thm:master} shows that the $1/\sqrt n$ scale is sharp in both
directions. We then prove the theorem that explains why one cannot do better by enumerating gate
patterns: any certificate required to hold over all patterns of a fixed density inherits the
worst-case conditioning, which is exponentially small by the product rate of \Cref{thm:product}.

\subsection{The Lipschitz certificate}
\begin{theorem}[Uniform robustness certificate]\label{thm:certificate}
Under the hypotheses of \Cref{thm:master} (fragility not required; only the Gaussian weights and the
wide regime), let
\[
  L_{\mathrm{net}}=\|w\|\prod_{\ell=1}^L\bigl(2\sqrt{\kappa(c_1)}\bigr)
\]
with $c_1$ the gate-density upper bound of \Cref{ass:nondeg}. Then w.h.p.\ $F$ is
$L_{\mathrm{net}}$-Lipschitz on the sphere, and consequently, for every input $x$, no perturbation of
Euclidean norm less than $|F(x)|/L_{\mathrm{net}}$ changes the sign of $F$. This certified radius holds
for \emph{all} inputs with no exceptional set.
\end{theorem}

\begin{proof}
For any $x,x'$ the difference $F(x')-F(x)$ is bounded by the operator norm of the Jacobian integrated
along the segment; on each cell of the arrangement (\Cref{lem:arrangement}) the Jacobian is a fixed
product of gated blocks, whose operator norm is at most
$\|w\|\prod_\ell\max_{|T_\ell|,|S_\ell|\le c_1 n_\ell}\|P_{T_\ell}W_\ell P_{S_\ell}\|_\op\le
\|w\|\prod_\ell(2\sqrt{\kappa(c_1)}+o(1))$ by \Cref{prop:A2}, uniformly over cells by the union bound
(the count is $(C\bar n)^{dL}$ and each operator-norm bound fails with probability $e^{-cn}$). Since
the gate densities are at most $c_1$ everywhere by \Cref{ass:nondeg}, this bounds the Lipschitz
constant on every cell, hence on the sphere. A perturbation smaller than $|F(x)|/L_{\mathrm{net}}$
changes $F$ by less than $|F(x)|$, so cannot change its sign.
\end{proof}

The subgaussian version, with an absolute constant $C(K)$ in place of $2\sqrt{\kappa(c_1)}$, holds for
any i.i.d.\ subgaussian weights of parameter $K$, by the subgaussian operator-norm bracket of
\Cref{app:brackets}; this is what makes the certificate meaningful for non-Gaussian and lightly-trained
weights.

\begin{corollary}[Sharpness of the $1/\sqrt n$ scale]\label{cor:sharp}
Under the hypotheses of \Cref{thm:master} with a random readout, for a fixed input $x$ (equivalently, for
all but a vanishing fraction of the sphere, excluding a neighborhood of the decision boundary where
$|F(x)|$ is atypically small), w.h.p.\ the certified radius and the fragility radius bracket each other:
\[
  \frac{c\,\|x\|}{\sqrt{n_L}\,L_{\mathrm{net}}}\ \le\ \bigl(\text{largest certified radius at }x\bigr)\ \le\
  \bigl(\text{attack radius at }x\bigr)\ \le\ \frac{C\,\|x\|\,\sqrt{d\log\bar n}}{\sqrt{n_L}} .
\]
Both scale as $\|x\|/\sqrt{n_L}$ up to the constants $L_{\mathrm{net}}=(2\sqrt{\kappa})^L$ and
$\prod\sqrt{\kcheck}$; the gap between them is $L_{\mathrm{net}}\big/\prod_\ell\sqrt{\kcheck(\alpha_\ell)}$,
a product of the two-sided gate constants and no power of $n$.
\end{corollary}

\begin{proof}
The lower bound is \Cref{thm:certificate} with $|F(x)|\ge c\|x\|/\sqrt{n_L}$, which holds w.h.p.\ at a
fixed $x$ because $F(x)\sim N(0,\|h_L\|^2/n_L)$ is not atypically small (anticoncentration of a
Gaussian); the upper bound is \Cref{thm:master}. The ratio of the two radii is
$L_{\mathrm{net}}\cdot(\text{attack constant})=(2\sqrt\kappa)^L\prod\sqrt{\kcheck}^{-1}$ up to
logarithms, independent of $n$.
\end{proof}

\subsection{The certification gap}
The certificate of \Cref{thm:certificate} is loose by the factor
$L_{\mathrm{net}}/\prod\sqrt{\kcheck}$ because it uses the worst-case operator norm at every step. One
might hope to sharpen it by certifying each gate pattern separately, using the actual conditioning of
the realized Jacobian rather than the worst pattern's. The product rate shows this cannot help beyond a
point: a certificate valid over all density-$\alpha$ patterns must survive the worst-conditioned one,
and that one is exponentially singular.

\begin{theorem}[Certification gap]\label{thm:cert-gap}
Consider certifying a depth-$k$ gated linear network $x\mapsto D_kW_k\cdots D_1W_1x$ with Gaussian
weights, by any method whose guarantee must hold uniformly over all gate tuples
$(D_1,\dots,D_k)$ of density $\alpha$. The smallest conditioning such a method must tolerate is the
least singular value of the worst pattern, $M_N^{(k)}(\alpha)$, which by \Cref{thm:product} is at most
$e^{-2H(\alpha)N/\dF+o(N)}$ w.h.p., and equal to $e^{-(k+1)H(\alpha)N/\dF+o(N)}$ at the sharp
(numerically determined) product rate. Consequently any pattern-uniform margin- or inversion-based
certificate degrades relative to the typical-pattern conditioning $N^{-(k+1)/2}$ by a factor at least
$e^{2H(\alpha)N/\dF+o(N)}$, and by $e^{(k+1)H(\alpha)N/\dF+o(N)}$ at the sharp rate: the union that a
uniform certificate performs is dominated by the single worst pattern, whose conditioning sets the rate.
\end{theorem}

\begin{proof}
A certificate uniform over the $\binom N{\alpha N}^{k+1}=e^{(k+1)H(\alpha)N+o(N)}$ patterns must in
particular be valid for the pattern minimizing $\smin$, whose value is $M_N^{(k)}(\alpha)$; by
\Cref{thm:product} this is $\le e^{-2H(\alpha)N/\dF+o(N)}$ w.h.p.\ (the proved achievability of
\Cref{prop:prod-lower}), and $=e^{-(k+1)H(\alpha)N/\dF+o(N)}$ at the sharp rate. A margin certificate
scales the certified radius by the least singular value of the map it inverts, so its worst-case radius
is smaller than the typical-pattern radius (governed by the fixed-support hard edge $N^{-(k+1)/2}$,
\Cref{prop:A4} and its least-sv companion) by the ratio $N^{-(k+1)/2}/M_N^{(k)}$, which is at least
$e^{2H(\alpha)N/\dF+o(N)}$ and equals $e^{(k+1)H(\alpha)N/\dF+o(N)}$ at the sharp rate. Either way it is
exponential in $N$.
\end{proof}

\begin{remark}[Interpretation]\label{rem:cert-interp}
\Cref{thm:cert-gap} is the precise sense in which the adversarial and robust pictures are two sides of
one random-matrix law. Fragility is the statement that the \emph{realized} pattern already conditions
the network badly enough for a gradient step to work; the certification gap is the statement that no
pattern-agnostic certificate can be sharp, because the \emph{worst} pattern is exponentially worse than
the realized one, at a rate that is at least $2H(\alpha)/\dF$ unconditionally and $(k+1)H(\alpha)/\dF$ at
the sharp product rate of \Cref{thm:product}. The only way
to beat the gap is to certify at the realized pattern --- i.e.\ to use the actual input, which is what a
per-input Lipschitz-of-the-realized-Jacobian certificate does; that recovers the $1/\sqrt n$ scale of
\Cref{cor:sharp} but is not uniform over patterns, and hence not robust to a gate flip, which is
exactly what an adversary induces.
\end{remark}

\subsection{Certified global robustness of sparse gates}
When the gate density is small enough that even the optimized product norm contracts, the network is
robust everywhere, with no small-perturbation attack possible.

\begin{theorem}[Sparse-gate robustness]\label{thm:sparse-robust}
Suppose \Cref{ass:nondeg} holds with $c_1<\alpha^*=0.042672$ (the root of $\kappa(\alpha^*)=1/4$). Then
w.h.p.\ the network Lipschitz constant $L_{\mathrm{net}}=(2\sqrt{\kappa(c_1)})^L\|w\|$ tends to $0$ as
$L\to\infty$, and for $L\ge L_0(c_1)$ the network computes a function whose oscillation on the sphere is
$o(1)$: no input has a certified radius smaller than $\Omega(1)$ relative to the output scale, so the
$1/\sqrt n$ fragility of \Cref{thm:master} is absent. Concretely, every perturbation of norm below
$|F(x)|/(2\sqrt{\kappa(c_1)})^L$ is harmless, and this radius grows with depth.
\end{theorem}

\begin{proof}
By \Cref{prop:contract}, $\kappa(c_1)<1/4$ makes $(2\sqrt{\kappa(c_1)})^L\to0$, so
$L_{\mathrm{net}}\to0$ and \Cref{thm:certificate} gives a certified radius $|F(x)|/L_{\mathrm{net}}$
that grows geometrically in $L$. The oscillation bound is the same operator-norm estimate applied to
$F(x)-F(x')$ over the whole sphere.
\end{proof}

This is the positive counterpart of the fragility theorems: sparsity of activation, quantified by the
gate density falling below $\alpha^*$, certifies global robustness, and the threshold is exactly the
contraction threshold of the optimized submatrix product. It is consistent with the converse of
\cite{bartlettBC21}, that networks deep enough to be near-constant are not fragile, and sharpens it to a
density condition with an explicit constant.
